% Build in this directory: pdflatex -interaction=nonstopmode -halt-on-error class1.tex
% Syllabus topics follow PT26w/lectureNote.tex.
\documentclass[a4paper,10pt,pdflatex,ja=standard,jafont=ipaex-type1,
  magstyle=real,scale=1]{bxjsarticle}
\usepackage{lmodern}
\usepackage{amsmath,amssymb,mathtools}
\usepackage{geometry}
\geometry{reset,left=22mm,right=22mm,top=20mm,bottom=17mm,
  headheight=13pt,headsep=6mm,footskip=10mm}
\usepackage{xcolor}
\usepackage{enumitem}
\usepackage{array,booktabs,tabularx}
\usepackage{fancyhdr}
\usepackage{hyperref}

\definecolor{ink}{HTML}{193A5C}
\definecolor{accent}{HTML}{237D82}
\definecolor{soft}{HTML}{EAF3F3}
\definecolor{muted}{HTML}{52616B}
\newcommand{\coursename}{確率論と統計学 I}
\hypersetup{unicode=true,colorlinks=true,urlcolor=accent,linkcolor=ink,
  pdfauthor={中島秀太},pdftitle={第1回：シラバスと確率空間}}
\pagestyle{fancy}
\fancyhf{}
\fancyhead[L]{\small\color{ink}\coursename{}・授業要約と演習}
\fancyhead[R]{\small\color{ink}第1回}
\fancyfoot[C]{\small\color{ink}\thepage/2}
\renewcommand{\headrulewidth}{0.3pt}
\setlength{\parindent}{0pt}
\setlength{\parskip}{2pt}
\setlength{\abovedisplayskip}{5pt}
\setlength{\belowdisplayskip}{5pt}
\setlength{\abovedisplayshortskip}{3pt}
\setlength{\belowdisplayshortskip}{4pt}
\setlist[itemize]{leftmargin=1.8em,itemsep=1pt,topsep=2pt,parsep=0pt}
\setlist[enumerate]{leftmargin=2em,itemsep=3pt,topsep=3pt,parsep=0pt}
\newcommand{\topic}[1]{\par\vspace{5pt}
  {\large\sffamily\mdseries\color{accent}#1}\par\vspace{2pt}}
\newcommand{\handouttitle}[1]{%
  \begin{center}
    {\Large\color{ink}#1}\par\vspace{3mm}
    {\color{accent}\rule{0.88\linewidth}{0.7pt}}
  \end{center}\vspace{-2mm}}
\newcommand{\keybox}[1]{\par\vspace{3pt}\begingroup
  \setlength{\fboxsep}{5pt}
  \colorbox{soft}{\parbox{\dimexpr\linewidth-2\fboxsep\relax}{#1}}%
  \endgroup\par\vspace{2pt}}
\newcommand{\examplelabel}[1]{{\sffamily\mdseries\color{ink}#1}}
\newcommand{\F}{\mathcal F}
\newcommand{\R}{\mathbb R}

\begin{document}
\handouttitle{シラバス}
{\large\color{ink}\coursename}\hfill
{\small 2026年度冬学期\quad 講義担当：中島 秀太}

\topic{授業の概要}
確率空間・確率変数・確率分布を出発点に、ルベーグ積分に基づく確率論を学ぶ。

\topic{到達目標}
\begin{itemize}
  \item 標本空間・事象・確率測度を区別し、具体的な確率モデルを構成できる。
  \item 確率変数の分布・期待値・独立性を説明し、基本的な計算ができる。
  \item 確率変数列の収束の違いを理解する。
\end{itemize}

\topic{授業計画（予定）}
\begingroup
\renewcommand{\arraystretch}{1.23}
\begin{tabularx}{\linewidth}{@{}>{\centering\arraybackslash}p{12mm}X@{}}
\toprule
回 & 内容 \\
\midrule
1 & 離散確率空間と一般の確率空間、$\sigma$-加法族 \\
2 & 確率変数と確率分布、分布関数 \\
3 & 期待値とルベーグ積分\\
4 & 独立性と共分散 \\
5 & 代表的な確率分布：二項・幾何・ポアソン・正規分布など \\
6 & 二項分布の正規近似 \\
7 & 特性関数とフーリエ解析、分布の一意性 \\
8 & 確率変数列の収束：概収束・確率収束 \\
9 & 分布の弱収束、極限定理 \\
10 & 大数の法則と一般の中心極限定理 \\
11 & 条件付き期待値、情報に関する平均 \\
12 & Kolmogorovの拡張定理、無限回の試行の確率空間 \\
13 & 総合演習と確率モデルへの応用 \\
\bottomrule
\end{tabularx}
\endgroup
\smallskip
{\small\color{muted}授業の進度に応じて、内容や順序を変更することがある。}

\topic{履修に必要な知識・学習の進め方}
集合と写像、数列・級数、微分積分、およびルベーグ積分の基礎を前提とする。
各回の定義を具体例で確かめ、講義後に証明と計算を復習する。
レポート問題の解答は、授業4回に一度、紙で提出する。
レポートでは、結論だけでなく、用いた仮定と理由を自分の言葉で記述する。

\topic{教材}
各回の配布資料を用いる。
関連する教科書：熊谷 隆『確率論』（新しい解析学の流れ）。

\topic{評価}
期末テスト 60\% + 授業レポート 20\% + 演習 20\%

総合得点が60\%以上で合格とし、合格者の成績は相対評価で決定する。

\newpage
\fontsize{10pt}{13.5pt}\selectfont
\setlength{\abovedisplayskip}{4pt}
\setlength{\belowdisplayskip}{4pt}
\handouttitle{第1回\quad 確率空間}

\topic{離散確率空間}
$\Omega$ を空でない可算集合（有限集合を含む）とする。
各結果 $\omega\in\Omega$ に重み $p_\omega$ を与え、
$$
  p_\omega\in[0,1]\quad(\omega\in\Omega),
  \qquad \sum_{\omega\in\Omega}p_\omega=1
$$
を仮定する。$\Omega$ を標本空間、その部分集合 $A\subseteq\Omega$ を事象と呼ぶ。
$$
{\ P(A)=\sum_{\omega\in A}p_\omega\ },
  \qquad \F=2^\Omega=\{A\mid A\subseteq\Omega\}.
$$
このとき $(\Omega,2^\Omega,P)$ は離散確率空間であり、
$p_\omega=P(\{\omega\})$ である。

\examplelabel{例：公平なコインを1回投げる。}
$H$ を表、$T$ を裏とすると、
$$
  \Omega=\{H,T\},\quad p_H=p_T=\frac12,\quad
  \F=\{\varnothing,\{H\},\{T\},\Omega\},\quad P(\{H\})=\frac12.
$$

\topic{一般の確率空間の公理}
$\Omega\ne\varnothing$ は集合、$\F\subseteq2^\Omega$ は
次を満たす $\sigma$-加法族とする。
$$
  \Omega\in\F,\qquad
  A\in\F\ \Longrightarrow\ A^c:=\Omega\setminus A\in\F,\qquad
  (A_n\in\F\ \forall n\ge1)\ \Longrightarrow\
  \bigcup_{n=1}^\infty A_n\in\F.
$$
次をみたす写像 $P:\F\to[0,1]$ は確率測度と呼ばれる:
$$
  P(\Omega)=1,\qquad
  P\!\left(\bigcup_{n=1}^\infty A_n\right)
    =\sum_{n=1}^\infty P(A_n)
  \quad\left(A_n\in\F,\ A_n\cap A_m=\varnothing\ (n\ne m)\right).
$$
三つ組 $(\Omega,\F,P)$ を確率空間と呼ぶ。公理から $P(\varnothing)=0$、$P(A^c)=1-P(A)$ が従う。

\examplelabel{例2：$[0,1]$ 上の一様分布。}
$$
  \Omega=[0,1],\qquad \F=\mathcal B([0,1]),\qquad
  P(A)=\lambda(A),\quad P([a,b])=b-a\quad(0\le a\le b\le1).
$$
$\mathcal B([0,1])$ を$[0,1]$に含まれる開集合から生成される最小の $\sigma$-加法族
（ボレル集合族）、$\lambda$ はルベーグ測度（長さ）である。


\topic{演習問題}
\begin{enumerate}[label=\arabic*.]
  \item 公平なサイコロを独立に2回投げる試行の離散確率空間 $(\Omega,\F,P)$ を構成し、
  出た目の合計が3となる確率を求めよ。
  \item $(\Omega,\F,P)$ を確率空間とする。任意の $A,B\in\F$ に対して、
  $A\cap B\in\F$ および $P(A\cup B)=P(A)+P(B)-P(A\cap B)$ を示せよ。
\end{enumerate}

\topic{レポート問題}
\begin{enumerate}[label=\arabic*.]
  \item 晴れの確率が30\%、くもりが50\%、雨が20\%となる確率空間を構成せよ。（ただし$\mathbb P$を具体的に構成する必要はなく、$p_{\omega}$のみでよい）
  \item 一般の確率空間で、確率を測れる集合を $\F$ に制限する理由を説明せよ。
  \item 第2問と関連して、自然に求める性質を保つために $\F$をすべての$\Omega$の部分集合 とできない例を一つ挙げ、
  どの性質が両立しないのか述べよ。
\end{enumerate}
\keybox{第2問はAIと相談してもよいが、その場合はどのような議論をしたかも書くこと。ただし、内容を自分で理解したうえで、
自分の言葉で説明すること。}
\end{document}
